\clearpage
\VerdierPhysicalPageAnchor{0124}%
\begingroup
\small
\setlength{\emergencystretch}{1.5em}%
\setlength{\abovedisplayskip}{4pt plus 2pt minus 2pt}%
\setlength{\belowdisplayskip}{4pt plus 2pt minus 2pt}%

\VerdierSemanticAnchor{ch2-1-3-2-triangle-morphism-conditions}
\noindent tels que
$(\operatorname{id}_{X_1},u_1,m'_1)$ et
$(u_3,\operatorname{id}_{X_3},m'_3)$ soient des morphismes de triangles et
tels que le triangle :

\VerdierSemanticAnchor{ch2-1-3-2-connecting-triangle}
\[
  Z'_3\xrightarrow{m'_1}Z'_2\xrightarrow{m'_3}Z'_1
  \xrightarrow{\,v'_3[1]w'_1\,}Z'_3[1]
\]
\noindent soit distingué.

\VerdierSemanticAnchor{ch2-1-3-2-comparison-isomorphisms}
De plus ((1.2.4) et (1.2.14)), il existe des isomorphismes
$i_j\colon Z'_j\longrightarrow Z_j$, $1\leq j\leq 3$, tels que
$(\operatorname{id},\operatorname{id},i_j)$ soient des isomorphismes de
triangles. Il suffit alors de poser :

\VerdierSemanticAnchor{ch2-1-3-2-transported-morphisms}
\[
  m_1=i_2m'_1i_3^{-1}\colon Z_3\longrightarrow Z_2,
  \qquad
  m_3=i_1m'_3i_2^{-1}\colon Z_2\longrightarrow Z_1.
\]

\VerdierSemanticAnchor{ch2-1-3-2-proof-complete}
Les morphismes $m_1$ et $m_3$ possèdent les propriétés demandées par
l'axiome TRIV.

\VerdierSemanticAnchor{ch2-1-3-3}
\medskip
\noindent\textbf{1.3.3.}\quad
Soit $C$ une catégorie \emph{additive}. À tout objet $X$ de $C$, on associe
le complexe de $C$ dont tous les composants sont nuls sauf le composant de
degré zéro qui est égal à $X$. On définit ainsi un foncteur de $C$ dans
$\operatorname{comp}(C)$, et en composant avec le foncteur canonique
$\operatorname{comp}(C)\longrightarrow K(C)$, un foncteur de $C$ dans
$K(C)$.

\VerdierSemanticAnchor{ch2-1-3-3-full-embeddings}
Ces foncteurs sont pleinement fidèles et injectifs sur les ensembles
d'objets, et réalisent donc $C$ comme sous-catégorie pleine de
$\operatorname{comp}(C)$ et $K(C)$. Par la suite, la catégorie $C$ sera
toujours considérée comme sous-catégorie pleine de $\operatorname{comp}(C)$
et $K(C)$ par l'intermédiaire de ces foncteurs.

\VerdierSemanticAnchor{ch2-1-3-4}
\medskip
\noindent\textbf{1.3.4.}\quad
Soit $C$ une catégorie \emph{abélienne}. Pour tout complexe $X$ de $C$, on
désigne par $H^0(X)$ l'objet de $C$ :

\VerdierSemanticAnchor{ch2-1-3-4-1}
\[
  H^0(X)=\operatorname{Ker}(d_X^0)/\operatorname{Im}(d_X^{-1}).
  \tag{1.3.4.1}
\]

\VerdierSemanticAnchor{ch2-1-3-4-induced-morphism}
Pour tout morphisme de complexes $X\xrightarrow{f}Y$, on désigne par
\[
  H^0(f)\colon H^0(X)\longrightarrow H^0(Y)
\]
\noindent le morphisme induit par $f$. On définit ainsi un foncteur :

\VerdierSemanticAnchor{ch2-1-3-4-comp-functor}
\[
  H^0\colon\operatorname{comp}(C)\longrightarrow C
\]
\noindent qu'on appelle le \emph{zéro-ième objet de cohomologie}. Ce foncteur
transforme les morphismes homotopes en morphismes égaux et, par suite,
fournit un foncteur encore noté :

\VerdierSemanticAnchor{ch2-1-3-4-k-functor}
\[
  H^0\colon K(C)\longrightarrow C.
\]
\endgroup
